\section{Part IV: Funding and Unit Economics of Scientific Work}

\subsection{27. The full cost of a research project}

A research budget includes salary and benefits, equipment or cloud use, consumables, collaborators, travel, participant costs, computing/storage, safety, administration, indirect facilities costs, audit and reporting, and a contingency appropriate to uncertainty. It must also include \textbf{time when a funded project cannot be charged to the grant}, and resources required to continue or terminate responsibly when funding ends.

Classify direct, indirect, restricted, and unrestricted costs. The institute can win a ₹50-lakh award and lose unrestricted cash if reimbursements are slow, overhead is underrecovered, or matching requirements are overlooked.

\subsection{28. Grants are not unrestricted revenue}

Grants may reimburse only eligible expenditures, prohibit certain equipment or staff uses, require milestones, and claw back unspent or misused amounts. Some are restricted contributions; some are exchange arrangements; accounting and tax treatment depends on the agreement and governing standards. A grant award announcement is not spendable cash.

Before applying, inspect eligible host status, principal-investigator eligibility, co-financing, overhead recovery, reporting schedules, release milestones, IP obligations, open-access rules, publication and data terms, audit rights, termination, and post-award liability. Public programs such as India's ANRF and BIRAC have distinct schemes and changing call windows; \textbf{never budget unawarded grants as committed financing} [R2-R3].

\subsection{29. Corporate sponsorship versus independent inquiry}

A client may pay for exclusive development and results; a research sponsor may support a question while permitting publication; a philanthropic donor may fund public scholarship without controlling conclusions. These are economically and ethically different contracts. Price and govern them accordingly.

If the sponsor can veto unfavorable findings, require deletion of inconvenient data, or manipulate authorship, the institution's independence is at risk. Narrow confidentiality and patent-filing delays may be legitimate; permanent sponsor control of scientific interpretation is a different bargain. State the publication review period, confidentiality, background IP, foreground IP, licensed fields, and termination rights precisely.

\subsection{30. Overhead recovery and the missing cost problem}

Suppose a research project budgets ₹30 lakh of direct costs and consumes ₹9 lakh of laboratories, accounting, insurance, safety, administration, and technology overhead. If the award permits only ₹3 lakh overhead, the institute must explicitly approve an additional ₹6 lakh subsidy. If project management calls the award ``₹33 lakh revenue'' and ignores the shortfall, the research center can grow grant volume while draining the parent organization.

Treat a permitted indirect-cost rate as a contract attribute-not proof that it reflects actual indirect cost. Create a separate ledger for full economic cost, sponsor reimbursement, internal subsidy, and cash timing.

\subsection{31. The research funding ladder}

At the smallest scale: founder-supported thought and reproducible open work. Then: cost-shared pilot investigations with clients or universities, paid research-service engagements with preserved reusable rights where lawful, competitive grants once eligible, sponsored chairs or facilities, donor-backed programs, and eventually an endowment for work requiring patient unrestricted time. These are options, not an automatic sequence.

Grant eligibility is frequently institution- and researcher-specific. A newly registered private course company may not qualify for programs available to recognized academic or R\&D institutions. The founder should build a \textbf{grant eligibility map}, not assume government research funds will finance an ambition simply because it uses the word research.

\subsection{32. Endowments and patient research capital}

[Model] A research group with annual recurring requirements of ₹1.8 crore would require a ₹45-crore corpus at an assumed 4\% annual spending rule, before fees, restrictions, inflation effects, and market volatility. The 4\% is an illustrative policy assumption, not a guaranteed return. An endowment invested for long horizons may experience losses; spending must be prudent and legally consistent with donor restrictions.

Early-stage researchers are often better financed by finite grants or transparent operating subsidy. A large permanently endowed chair is a late funding architecture, not a condition for doing the 1st serious experiment.

\subsection{33. A 12-month research-lab budget}

[Model: software/methods laboratory, not a real Bayes quote]


\begin{table}[t]
\centering
\small
\caption{12-Month Research-Lab Budget Model (Software/Methods Laboratory)}
\label{tab:ch33_lab_budget}
\begin{tabularx}{\columnwidth}{X r}
\toprule
\textbf{Item} & \textbf{Annual Cash Cost (\textnormal{₹} lakh)} \\
\midrule
Research lead, fractional or contracted & 12 \\
Research engineer & 15 \\
Research assistance and annotation & 6 \\
Cloud compute, storage, software & 8 \\
External review and collaboration & 3 \\
Publication, security, data, travel & 3 \\
Administrative and facilities allocation & 5 \\
Contingency & 4 \\
\midrule
\textbf{Full planned annual cash envelope} & \textbf{56} \\
\bottomrule
\end{tabularx}
\end{table}


The table implies ₹4.67 lakh average monthly expense, but grant receipts and equipment purchases may be uneven. It excludes any unpriced founder involvement and may be insufficient for particular senior talent or difficult compute workloads. Do not hire against this budget without secured financing and a runway policy.

\subsection{34. Cash timing: the reimbursable grant trap}

[Model] A ₹30-lakh approved project reimburses 70\% of eligible expenditure quarterly after reports. The institute must pay ₹10 lakh in the 1st quarter while receiving nothing until verification. Even if the contract ultimately funds ₹21 lakh, the temporary working-capital requirement may be large and the remaining ₹9 lakh plus overhead must be funded elsewhere.

Build a monthly restricted-cash forecast, drawdown schedule, evidence submission checklist, and recovery reserve. In stress cases assume delayed releases, disallowed costs, exchange movements, and delayed recruitment. Project profitability cannot substitute for cash available to make payroll.

\subsection{35. The research investment committee}

Before committing beyond a small founder-approved threshold, convene at least 1 scientifically qualified independent reviewer and 1 financially responsible decision maker. Ask: Is the question important? Is there a credible approach and replicability plan? Who leads? What are ethical, contractual, or publication constraints? What is fully loaded cost? What are the staged learning gates? What must be true for the next allocation? Who can halt unsafe or fraudulent work?

Record conflicts and dissent. The founder may allocate capital but should not be the sole judge of claims outside their expertise.
